\section{Discusión}

\subsection{Fuerza restauradora y carácter armónico simple}
En el péndulo simple, la componente tangencial del peso actúa como fuerza restauradora: $F_t=-mg\sin\theta$. Aplicando la segunda ley de Newton para la rotación en torno al eje de suspensión,

\begin{equation}
mL\,\ddot{\theta}=-mg\sin\theta,
\end{equation}

de donde la masa se cancela, quedando

\begin{equation}
\ddot{\theta}+\frac{g}{L}\sin\theta=0.
\end{equation}

Para ángulos pequeños, $\sin\theta\approx\theta$, por lo que la ecuación se reduce a la forma canónica de un oscilador armónico simple,

\begin{equation}
\ddot{\theta}+\omega^2\theta=0, \qquad \omega=\sqrt{\frac{g}{L}},
\end{equation}

cuya solución es $\theta(t)=\theta_0\cos(\omega t+\varphi)$, con período $T=2\pi\sqrt{L/g}$. Este resultado es consistente con dos observaciones directas de los datos. La cancelación de $m$ en la ecuación de movimiento concuerda con que los períodos medidos para las tres esferas (Tabla~\ref{tab:tabla_masas}: $2.08$, $2.05$ y $2.07\,\mathrm{s}$, con desviaciones de \SI{1.41}{\percent}, \SI{-0.05}{\percent} y \SI{0.92}{\percent} respecto al valor teórico) sean similares pese a las diferencias entre las masas. Asimismo, en la Tabla~\ref{tab:tabla_amplitudes} el período se mantiene en $1.49\,\mathrm{s}$ para $\theta_0=30^\circ$ y $20^\circ$, con una variación de \SI{1.34}{\percent} hasta $\theta_0=15^\circ$. Los valores experimentales obtenidos para las distintas amplitudes muestran variaciones pequeñas del período, lo cual es consistente con la aproximación de pequeñas oscilaciones en el intervalo estudiado.

\subsection{¿Puede considerarse nuestro péndulo un péndulo simple?}
El modelo de péndulo simple idealiza el sistema como una masa puntual suspendida de un hilo inextensible y sin masa. Los datos se aproximan razonablemente a este ideal: las desviaciones porcentuales del período respecto al valor teórico fueron menores al \SI{1.5}{\percent} en todos los casos de la Tabla~\ref{tab:tabla_masas}, y los datos de la Tabla~\ref{tab:tabla_longitudes} muestran una tendencia general decreciente compatible con $T\propto\sqrt{L}$ al reducir $L$ de $99.5\,\mathrm{cm}$ a $79.5\,\mathrm{cm}$. Sin embargo, se observa una desviación local entre algunas configuraciones. Esta puede asociarse a las limitaciones del montaje real respecto al modelo ideal, entre ellas la masa y elasticidad del cordel, el tamaño finito de la esfera, el rozamiento del sistema y las incertidumbres en la medición de la longitud y del período. En conjunto, los resultados indican que el montaje se comporta como un péndulo simple \emph{en buena aproximación}, al presentar pequeñas desviaciones respecto al modelo teórico en el estudio de la masa y pequeñas variaciones del período dentro del intervalo de amplitudes analizado, aunque no reproduce de forma exacta todas las condiciones del modelo ideal.

\subsection{Relación entre la amplitud y el tiempo en el régimen amortiguado}
El ajuste exponencial obtenido a partir de los 15 máximos de amplitud registrados en la Tabla~\ref{tab:tabla_amortiguado}, $y=(173.80\pm0.64)e^{(-0.00573\pm0.00027)x}$ con $R^2=0.9693$, muestra una buena concordancia con una relación de la forma $\theta_{\max}(t)=\theta_0e^{-\gamma t}$. El ajuste presenta una buena correspondencia con los datos experimentales, como indica $R^2=0.9693$ y se observa en la Fig.~\ref{fig:regresion_exponencial}. El comportamiento observado es compatible con el modelo de amortiguamiento viscoso lineal utilizado ($F_r=-bv$), sin que ello permita atribuir de manera exclusiva la disipación a este mecanismo.

\subsection{Cumplimiento de los objetivos de la práctica}
Los resultados obtenidos son consistentes con los objetivos planteados. No se observó una dependencia significativa del período con la masa (Tabla~\ref{tab:tabla_masas}); el período presentó una tendencia general a disminuir cuando disminuyó la longitud (Tabla~\ref{tab:tabla_longitudes}); y, para las amplitudes ensayadas, permaneció aproximadamente constante (Tabla~\ref{tab:tabla_amplitudes}). Asimismo, la amplitud disminuyó con el tiempo (Tabla~\ref{tab:tabla_amortiguado}) y presentó una buena concordancia con un modelo exponencial. El ajuste mostró un comportamiento compatible con el amortiguamiento viscoso lineal y produjo $b=(6.65\pm0.32)\times10^{-5}\,\mathrm{kg/s}$.
